%% verse_protrusion.tex
%% Microtype protrusion for verse numbers (edition-portable).
%% Requires, loaded BEFORE this file: fontspec (main font set), microtype, scripture.
%% Effect: a verse number at the start of a line hangs a tuned partial amount
%% into the left margin. Mid-line numbers are unaffected (microtype only
%% protrudes line-edge glyphs). Font-agnostic: uses OpenType superior figures
%% of the CURRENT body font; falls back to \textsuperscript if unavailable.
%%
%% WORKING MECHANISM (the bits that are easy to get wrong):
%%   1. Use [ context = verseprot ], NOT [ name = ... ]. The `context' key is
%%      what makes \microtypecontext{protrusion=verseprot} resolve to this list;
%%      `name' only defines a list and leaves the context undefined (microtype
%%      then warns "I cannot find a protrusion list ... Switching off").
%%   2. Font set is the full 5-axis wildcard { font = */*/*/*/* } so it matches
%%      BOTH the body font instance AND the separate font instance that fontspec
%%      spins up for VerticalPosition=Superior (e.g. EBGaramond(1)).
%%   3. Each entry is char = {left,right} (ONE brace, comma-separated). Left-only
%%      hang -> {350,0}. Two separate braces {350}{0} is wrong syntax.
%%   4. microtype does NOT expand a macro inside the value braces, so the factor
%%      is an inlined literal (350) in every entry, not \verseprot@factor.
%%   5. *** THE KEY FIX ***  microtype (luatex) keys protrusion by the codepoint
%%      of the glyph AS IT APPEARS IN THE NODE LIST, i.e. AFTER GSUB shaping.
%%      VerticalPosition=Superior runs the OpenType `sups' feature, which in
%%      EB Garamond substitutes each ASCII digit for the glyph at its Unicode
%%      SUPERSCRIPT codepoint:
%%          0->U+2070  1->U+00B9  2->U+00B2  3->U+00B3
%%          4..9 -> U+2074..U+2079
%%      Therefore the list MUST be keyed on those superscript codepoints, NOT on
%%      the ASCII digits 0-9. (Keying 0-9 = U+0030..U+0039 was the bug: those
%%      input glyphs never reach the line because sups replaces them, so the
%%      lpcode never bound to the rendered superior glyph -> ZERO protrusion,
%%      verified by 600 dpi pixel measurement and by dumping the substituted
%%      glyph names with fontTools: three -> uni00B3, etc.)
%%   6. microtype protrudes GLYPH nodes, not boxes. The substituted superior
%%      figures are real glyphs and protrude; \textsuperscript (a raised box)
%%      does not -- hence the OpenType path is required for the hang, with
%%      \textsuperscript only as the non-protruding fallback when the font lacks
%%      superior figures.
%% Verified by measurement (tests/verse_protrusion_visual_check.tex, 600 dpi):
%% with this setup the leading verse number's leftmost ink sits measurably LEFT
%% of the column rule, while mid-line numbers stay inline; feature ON vs OFF
%% now differs at the left edge.
%%
%% PORTABILITY NOTE: the superscript-codepoint mapping above is what EB Garamond
%% (and most fonts following the Unicode/AGL convention) emit for `sups'. A font
%% whose superior figures live at private-use or `.sups'-named glyph slots with
%% no Unicode value would need different keys; in that case microtype cannot
%% address them by codepoint and the \textsuperscript fallback (no protrusion)
%% is the safe outcome.
\makeatletter

% --- Tunable: partial left-protrusion factor (out of 1000 = one glyph width).
%     NOTE: microtype does NOT expand a macro inside the \SetProtrusion value
%     braces (it reads the token name literally -> "Missing number"). So the
%     factor MUST be an inlined literal in every digit entry below. Keep this
%     def in sync with the literals; it is documentation only.
\def\verseprot@factor{350}

% --- Manual override (an edition may force the feature off):
\newif\ifverseprotrusion  \verseprotrusiontrue

% --- microtype protrusion context: left-only hang on the SUPERSCRIPT glyphs
%     that the `sups' feature substitutes for the digits 0-9 (see header note
%     #5). Keyed on the substituted Unicode codepoints, NOT on ASCII 0-9.
%     Activated ONLY inside \verseprotrudenum via \microtypecontext, so no
%     other text is affected.
\SetProtrusion
  [ context = verseprot ]
  { font = */*/*/*/* }
  { "2070 = {350,0} ,  % superscript 0
    "00B9 = {350,0} ,  % superscript 1
    "00B2 = {350,0} ,  % superscript 2
    "00B3 = {350,0} ,  % superscript 3
    "2074 = {350,0} ,  % superscript 4
    "2075 = {350,0} ,  % superscript 5
    "2076 = {350,0} ,  % superscript 6
    "2077 = {350,0} ,  % superscript 7
    "2078 = {350,0} ,  % superscript 8
    "2079 = {350,0} }  % superscript 9

% --- Feature detection (by glyph-width measurement) + final decision.
%     Done at begin-document when the real body font is active.
\newif\ifverseprot@use
\AtBeginDocument{%
  \setbox0=\hbox{{\addfontfeature{VerticalPosition=Superior}0}}%
  \setbox2=\hbox{0}%
  \ifdim\wd0=\wd2 % no substitution happened -> font lacks superior figures
    \verseprot@usefalse
  \else
    \ifverseprotrusion \verseprot@usetrue \else \verseprot@usefalse \fi
  \fi
}

% --- The formatter wired into scripture's verse/format.
\NewDocumentCommand\verseprotrudenum{m}{%
  \ifverseprot@use
    {\addfontfeature{VerticalPosition=Superior}%
     \microtypecontext{protrusion=verseprot}#1}%
  \else
    \textsuperscript{#1}%
  \fi
}

\makeatother
